% Part: first-order-logic
% Chapter: introduction
% Section: soundness-completeness

\documentclass[../../../include/open-logic-section]{subfiles}

\begin{document}

\olfileid{fol}{int}{scp}

\olsection{निर्दोषता आणि संपूर्णता}

प्रथम-क्रम तर्कशास्त्रासाठीच्या !!{derivation} पद्धतींचाही आपण परिचय करून
घेऊ. तर्कशास्त्रज्ञांनी अनेक !!{derivation} पद्धती विकसित केल्या आहेत,
पण त्या सर्व !!{sentence}s दरम्यान तोच !!{derivability} संबंध परिभाषित
करतात. विशिष्ट आणि काटेकोरपणे परिभाषित प्रकारची !!a{derivation} असेल,
तर~$\Gamma$ हे~$!A$ \emph{!!{derive}s} करते, $\Gamma \Proves !A$, असे
आपण म्हणतो. !!^{derivation}s म्हणजे नेहमी चिन्हांच्या सांत मांडण्या
असतात---कदाचित !!{sentence}s ची यादी किंवा त्याहून गुंतागुंतीची काही
संरचना. एखाद्या संच~$\Gamma$ पासून !!a{sentence} तार्किकरीत्या निष्पन्न
होते का हे ठरवण्याचे साधन पुरवणे हे !!{derivation} पद्धतींचे उद्दिष्ट
आहे. ते उद्दिष्ट साध्य करण्यासाठी $\Gamma \Entails !A$ तेव्हा आणि केवळ
तेव्हाच $\Gamma \Proves !A$ असले पाहिजे.

$\Gamma \Proves !A$ पण $\Gamma \Entails !A$ नसेल, तर आपली
!!{derivation} पद्धत अतिप्रबल असेल, म्हणजे ती गरजेपेक्षा जास्त सिद्ध
करेल. $\Gamma \Proves !A$ असेल तर $\Gamma \Entails !A$ असते या
गुणधर्माला \emph{निर्दोषता} म्हणतात आणि कोणत्याही चांगल्या
!!{derivation} पद्धतीसाठी ही किमान आवश्यकता आहे. दुसरीकडे,
$\Gamma \Entails !A$ पण $\Gamma \Proves !A$ नसेल, तर आपली
!!{derivation} पद्धत अतिदुर्बल असेल, म्हणजे ती पुरेसे सिद्ध करणार नाही.
$\Gamma \Entails !A$ असेल तर $\Gamma \Proves !A$ असते या गुणधर्माला
\emph{संपूर्णता} म्हणतात. निर्दोषता सिद्ध करणे सहसा तुलनेने सोपे असते
(विगमनाधारित रीतीने परिभाषित केलेल्या !!{derivation}s च्या संरचनेवर
विगमन करून). संपूर्णता सिद्ध करणे अधिक कठीण असते.

निर्दोषता आणि संपूर्णता यांचे अनेक महत्त्वाचे परिणाम आहेत.
!!{sentence}s चा एखादा संच~$\Gamma$ हा व्याघात (जसे
$!A \land \lnot !A$) !!{derive}s करत असेल, तर त्याला \emph{विसंगत}
म्हणतात. विसंगत~$\Gamma$ ला कोणतेही प्रतिमान असू शकत नाही; तो
अपूर्ततायोग्य असतो. संपूर्णतेवरून याचा उलट निष्कर्ष मिळतो: विसंगत
नसलेल्या---किंवा, आपण म्हणू त्याप्रमाणे, \emph{सुसंगत}---प्रत्येक
$\Gamma$ ला प्रतिमान असते. खरे तर, हे संपूर्णतेशी सममूल्य आहे आणि
संपूर्णतेचे हेच रूप आपण प्रत्यक्षात सिद्ध करू. हा सखोल आणि कदाचित
आश्चर्यकारक परिणाम आहे: $\Gamma$ पासून $!A \land \lnot !A$ सिद्ध करता
येत नाही एवढ्यामुळेच~$\Gamma$ वर्णन करते तसे !!a{structure} अस्तित्वात
असण्याची हमी मिळते. म्हणून संपूर्णता पुढील प्रश्नाचे उत्तर देते:
!!{sentence}s च्या कोणत्या संचांना प्रतिमाने असतात? उत्तर: सर्व आणि फक्त
सुसंगत संचांना.

निर्दोषता आणि संपूर्णता प्रमेयांचे दोन महत्त्वाचे परिणाम आहेत: संहतता
प्रमेय आणि लोव्हेनहाइम--स्कोलेम प्रमेय. प्रतिमानांच्या उपपत्तीतील हे
महत्त्वाचे परिणाम आहेत आणि अनेक रोचक निष्कर्ष प्रस्थापित करण्यासाठी
त्यांचा उपयोग करता येतो. आपण आधीच दोन निष्कर्षांचा उल्लेख केला आहे:
प्रथम-क्रम तर्कशास्त्रात !!a{structure} चा !!{domain} सांत आहे किंवा तो
!!{nonenumerable} आहे असे व्यक्त करता येत नाही.

ऐतिहासिकदृष्ट्या हे सर्व---प्रथम-क्रम तर्कशास्त्राची विन्यासमीमांसा आणि
चिन्हार्थमीमांसा कशी परिभाषित करायची, चांगल्या !!{derivation} पद्धती
कशा परिभाषित करायच्या, त्या निर्दोष आणि संपूर्ण आहेत हे कसे सिद्ध
करायचे, प्रथम-क्रम भाषांत काय व्यक्त करता येते आणि काय व्यक्त करता येत
नाही हे स्पष्ट करणे---शोधून योग्य रीतीने ठरवण्यास बराच काळ लागला. हे
कसे करायचे हे आता माहीत असले, तरी सर्व तपशील समजून घेणे अजूनही
गोंधळात टाकणारे आणि कंटाळवाणे ठरू शकते. पण ते महत्त्वाचेही आहे, कारण
प्रथम-क्रम तर्कशास्त्राच्या आकारिक भाषेसाठी येथे विकसित केलेल्या पद्धती
तर्कशास्त्र, संगणकशास्त्र आणि भाषाविज्ञान यांत सर्वत्र वापरल्या जातात.
म्हणून सर्व तपशीलांचा अभ्यास करण्याचे दीर्घकाळात फळ मिळते.

\end{document}
